\chapter{Estado del arte}
\label{ch:ch2_soa}

% [BORRADOR v1] Material de partida: docs/estado_del_arte_trng.md, estado_del_arte_cripto.md,
% plataforma_y_flujo.md, via_asic.md. Cada afirmación numérica de este capítulo lleva una
% referencia que se cotejó con su fuente al redactar la revisión;
% lo que la revisión no pudo consultar está marcado como tal en las notas de trabajo.

La revisión se organiza siguiendo la cadena que construye el TFM: cómo
se convierten los osciladores en anillo en bits aleatorios, cómo se
modela y se acota su entropía, cómo se mide su jitter desde dentro del
dispositivo, qué exigen las dos normas que rigen el campo, qué primitivas
simétricas encajan en un motor con recursos limitados, qué ofrecen los
componentes comerciales y qué cambia cuando el diseño sale de la FPGA
hacia el silicio. La última sección expone el hueco que aborda este
trabajo, en la forma prudente que exigen las reglas de trabajo.

\section{Arquitecturas de TRNG basadas en osciladores en anillo}
\label{sec:soa_archs}

Un oscilador en anillo es una cadena impar de inversores cerrada sobre sí
misma; su periodo fluctúa porque cada transición se retrasa una cantidad
aleatoria, y esa fluctuación --- el jitter --- es la fuente de ruido de
toda una familia de generadores. Lo que distingue a las arquitecturas es
cómo se acumula y se muestrea el jitter~\cite{petura_2016_survey, lubicz_2024_recommendations}.

El \textbf{oscilador en anillo elemental} (ERO) muestrea un anillo con un
flip-flop cuyo reloj es un segundo anillo dividido por $K_D$; el divisor
fija el tiempo de acumulación. Tiene el modelo estocástico más
maduro~\cite{baudet_2011_security, killmann_2008_design, ma_2014_entropy} y
una medida embebida del jitter~\cite{fischer_2014_embedded}, a costa de
una tasa de salida muy baja ($K_D$ de $10^4$ a $10^5$: 80\,000 en
Spartan-6 y 135\,000 en Cyclone~V en la implementación de Petura, unos
430\,000 en el diseño de Fischer y Lubicz para 0,997~bit de entropía). El
generador \textbf{multianillo} de Sunar, Martin y
Stinson~\cite{sunar_2007_provably} combina mediante XOR muchos anillos
muestreados por un reloj de referencia; Wold y Tan~\cite{wold_2009_analysis}
añadieron un flip-flop por anillo antes de la XOR y publicaron 100~Mbit/s
con menos de 100 elementos lógicos. Su debilidad la pusieron de manifiesto
Bochard \emph{et al.}~\cite{bochard_2010_true}: parte de la aleatoriedad
aparente es \emph{pseudo}aleatoriedad debida a la interacción entre
anillos de frecuencias distintas, y la independencia que supone la
demostración no se cumple en una FPGA. El \textbf{muestreo coherente}
(COSO) usa dos anillos de periodo casi igual y cuenta el batido~\cite{kohlbrenner_2004_embedded,
valtchanov_2009_enhanced}; su modelo da la min-entropía directamente a
partir de la distribución de la cuenta~\cite{peetermans_2021_configurable},
y la calibración dinámica de anillos reconfigurables lo hace portable
entre dispositivos y emplazamientos (3,30~Mbit/s en Spartan-6,
4,65~Mbit/s en Spartan-7)~\cite{peetermans_2019_portable, peetermans_2021_configurable}.
El \textbf{anillo de efecto de transición} (TERO) cuenta cuántas veces
oscila una estructura biestable antes de estabilizarse~\cite{varchola_2010_new},
con un modelo completo que va de lo físico a lo
estocástico~\cite{bernard_2019_physical}; el \textbf{anillo
autotemporizado} (STR) hace circular varios eventos por un anillo
asíncrono para obtener fases separadas del orden del
jitter~\cite{cherkaoui_2013_very}, y alcanza 154~Mbit/s; y el
\textbf{ES-TRNG} de Yang \emph{et al.} usa la cadena de acarreo como
convertidor de tiempo a digital para codificar la posición del flanco con
10~LUT y 5~flip-flops~\cite{yang_2018_estrng, rozic_2015_highly}.

La Tabla~\ref{tab:petura} reproduce la comparación que hicieron Petura
\emph{et al.} de los núcleos conformes a AIS sobre un mismo
dispositivo~\cite{petura_2016_survey}. La columna de entropía es una
estimación estadística (el test T8 de AIS~31), no una cota del modelo; la
columna de viabilidad es la valoración cualitativa de los autores, en la
que COSO y TERO puntúan peor porque necesitan un ajuste manual por
dispositivo. No existe una tabla comparable para Artix-7; el proyecto
abierto OpenTRNG implementa ERO, multianillo y COSO en una Arty~A7 con el
mismo dado XC7A35T que la Basys~3, pero no publica cifras ni de entropía
ni de tasa de bits~\cite{opentrng_2025_docs}.

\begin{table}[htbp]
\centering
\caption{Núcleos TRNG compatibles con AIS en Xilinx Spartan-6 (45~nm),
según Petura \emph{et al.}~\cite{petura_2016_survey}, tabla~II. La
entropía es la estimación T8; la ``viabilidad'' es la clasificación de los
autores (cuanto más alta, más fácil).}
\label{tab:petura}
\begin{tabular}{@{}lrrrrl@{}}
\toprule
Núcleo & LUT/reg. & Potencia (mW) & Tasa (Mbit/s) & Entropía/bit & Viab. \\
\midrule
ERO  & 46/19   & 2,16  & 0,0042 & 0,999 & 5 \\
COSO & 18/3    & 1,22  & 0,54   & 0,999 & 1 \\
MURO & 521/131 & 54,72 & 2,57   & 0,999 & 4 \\
PLL  & 34/14   & 10,6  & 0,44   & 0,981 & 3 \\
TERO & 39/12   & 3,31  & 0,625  & 0,999 & 1 \\
STR  & 346/256 & 65,9  & 154    & 0,998 & 2 \\
\bottomrule
\end{tabular}
\end{table}

\section{Modelos estocásticos y cotas de entropía}
\label{sec:soa_models}

La exigencia de que un generador físico vaya acompañado de un
\emph{modelo estocástico} --- una descripción matemática de la fuente de
ruido a partir de la cual puede derivarse la entropía de la salida --- la
introdujeron Killmann y Schindler~\cite{killmann_2008_design, schindler_2002_evaluation}
y es obligatoria en AIS~20/31. Para los osciladores en anillo, el modelo
de referencia es el de Baudet \emph{et al.}~\cite{baudet_2011_security}:
la fase del anillo muestreado es un proceso de Wiener con deriva $\mu$ y
volatilidad $\sigma^2$, el bit muestreado es el semiperiodo en el que cae
la fase, y todo depende del factor de calidad $Q=\sigma^2\Delta t$
(varianza de fase acumulada entre muestras) y de la razón de frecuencias
$\nu=\mu\Delta t$. El modelo da la probabilidad del bit condicionada a la
fase anterior, la probabilidad de cualquier vector de salida y una cota
inferior de la tasa de entropía de Shannon,
$H \ge 1 - \frac{4}{\pi^2\ln 2}e^{-4\pi^2 Q}$, cuyo alcance exacto se
examina en el Capítulo~\ref{ch:ch3_entropy}. Petura \emph{et al.} la
reescriben en magnitudes de diseño, $H = 1 - \frac{4}{\pi^2\ln2}\exp(-\pi^2
\sigma_{th}^2 K T_2/T_1^3)$, y Fischer y Lubicz la invierten para obtener
el divisor $K_D$ a partir de un jitter medido~\cite{fischer_2014_embedded}.
Ma \emph{et al.} proponen un modelo alternativo que usa el grupo de
Lovaina~\cite{ma_2014_entropy}; Haddad \emph{et al.} cuestionan la
independencia mutua de las realizaciones del jitter en la que se apoyan
estas cancelaciones~\cite{haddad_2014_assumption}; y Saarinen advierte de
que la cota de Baudet ``nunca es inferior a 0,415, ni siquiera cuando $Q$
tiende a cero'' y solo es segura bajo hipótesis adicionales, y propone
una cota válida en todo el intervalo~\cite{saarinen_2021_entropy}.

Hay dos matices que importan en el diseño. Primero, no todo el jitter es
entropía. Fischer y Lubicz separan la varianza de fase acumulada en un
término de paseo aleatorio, lineal en el tiempo, producido por el ruido
térmico (blanco), y unos términos de parpadeo ($1/f^\beta$) cuya varianza
crece cuadráticamente y que dominan con tiempos de acumulación
largos~\cite{fischer_2014_embedded}; el jitter determinista y global que
procede de la alimentación y del sustrato se comporta como estos últimos
y puede ser conocido o inducido por un
adversario~\cite{fischer_2008_enhancing}. Solo puede computarse el
término lineal, y por eso no se puede hacer el divisor arbitrariamente
grande para compensar una fuente débil. Segundo, toda medida embebida es
\emph{diferencial}: dos anillos que comparten alimentación ven el mismo
jitter global, que se cancela en su fase relativa y deja la suma de sus
jitters locales. Lubicz y Sk\'orski formalizan este ``principio de
transferencia del jitter'' y muestran que con dos anillos solo se obtiene
la suma, mientras que con tres anillos y tres medidas por pares se
recuperan las varianzas individuales; su experimento muestra además que
la hipótesis de varianza lineal con el periodo falla en un factor dos en
uno de los casos~\cite{lubicz_2024_jitter}.

\section{Medida embebida del jitter}
\label{sec:soa_jitter}

Los osciloscopios no pueden resolver los picosegundos de jitter de un
anillo dentro de un chip sin un divisor externo~\cite{benea_2024_flicker},
así que los métodos útiles miden desde dentro. El \textbf{método del
contador} cuenta periodos de un anillo en una ventana fija del otro; la
varianza de la cuenta crece con el jitter relativo
acumulado~\cite{killmann_2008_design,
lubicz_2015_towards}, con un suelo de cuantificación de una cuenta y con
contaminación por ruido de parpadeo en ventanas largas. El método de
\textbf{Fischer y Lubicz}~\cite{fischer_2014_embedded} no necesita
divisor: en $K\approx10^4$ bloques de $N\approx100$ bits muestreados con
$K_D=1$, la fracción de pares $(b_j,b_{j+M})$ que difieren sigue a la
fase acumulada en $M$ periodos, y su varianza entre bloques, representada
frente a $M$ entre 200 y 1600, da el jitter por periodo con un error
inferior al 5\,\% en simulación; en una Cyclone~III midieron
$\sigma=5{,}01$~ps por periodo de 7,81~ns, y mostraron que enfriar el chip
reduce $\sigma$ --- lo que convierte esa misma medida en un test en
línea. El \textbf{muestreo coherente} proporciona el jitter acumulado a
partir de la varianza de la cuenta de batido sin hardware
adicional~\cite{peetermans_2021_configurable}. Más recientemente, Benea
\emph{et al.} caracterizan el jitter de los anillos con la
\textbf{varianza de Allan}, ajustando $\sigma_t^2(t)=a_0+a_1 t+a_2 t^2$
para separar los términos de cuantificación, térmico y de parpadeo, y lo
hacen tanto en una Artix-7 --- la familia de FPGA de este trabajo --- como
en un circuito FD-SOI de 28~nm, con los datos y los scripts
publicados~\cite{benea_2024_flicker}. Las cifras publicadas de jitter por
periodo se agrupan en 2--4~ps para periodos de 4--8~ns en Spartan-6 y
Cyclone~V~\cite{petura_2016_survey}, 5~ps en Cyclone~III con medida
embebida~\cite{fischer_2014_embedded}, y 2--2,5~ps en Cyclone~III y
Virtex-5 para anillos autotemporizados~\cite{cherkaoui_2013_very}; los
coeficientes de Benea \emph{et al.} para Artix-7 se publican en unidades
del contador, no en segundos, así que queda por medir una cifra absoluta
para la XC7A35T.

\section{Normas: AIS 20/31 v3.0 y NIST SP 800-90B}
\label{sec:soa_standards}

Dos documentos rigen cómo se evalúa un generador físico. La norma alemana
\textbf{AIS~20/31}, en su versión~3.0 de septiembre de 2024~\cite{peter_2024_ais31},
sustituye la exigencia de 2011 de 0,997~bit de entropía de Shannon por
bit interno~\cite{killmann_2011_ais31} por una alternativa en la clase
PTG.2: entropía de Shannon $\ge0{,}9998$ y/o min-entropía $\ge0{,}98$,
junto con $P(1)\in(0{,}493;\ 0{,}507)$; la política de transición del BSI
solo permite nuevas evaluaciones con la versión antigua hasta marzo de
2026~\cite{bsi_2025_transition}. Sus requisitos son tan estructurales
como numéricos: un modelo estocástico verificable es obligatorio; el test
en línea debe derivarse del modelo y fallar con alta probabilidad cuando
el parámetro de la fuente sale del conjunto admisible; un test de fallo
total debe detectar que la entropía cae a cero, y para ello se acepta
explícitamente la prueba de recuento de repeticiones de SP~800-90B; y la
batería de caja negra $T_{\mathrm{irn}}$ incluye ahora dos de los
predictores del NIST (MultiMMC y LZ78Y), una armonización
deliberada~\cite{peter_2024_ais31, schindler_2024_ecw}.

La norma estadounidense \textbf{SP~800-90B}~\cite{turan_2018_sp80090b}
estima la \emph{min-entropía} --- la magnitud que cuenta frente a un
adversario, $-\log_2$ de la probabilidad del valor más probable --- a
partir de al menos un
millón de muestras en bruto y de una matriz de reinicios de
1000$\times$1000, con el estimador del valor más común si se demuestra
que la fuente es i.i.d.\ y, en caso contrario, con el mínimo de diez
estimadores, cuatro de ellos predictores que intentan adivinar la muestra
siguiente. Aprueba dos tests de salud continuos y no exige ningún otro:
la prueba de recuento de repeticiones, con umbral
$C=1+\lceil-\log_2\alpha/H\rceil$, y la prueba de proporción adaptativa
sobre una ventana de 1024 muestras para fuentes binarias (512 para las no
binarias), $C=1+\mathrm{CRITBINOM}(W,2^{-H},1-\alpha)$, con
$\alpha=2^{-20}$ como valor recomendado. Enumera los componentes de
acondicionamiento cuya salida puede declarar entropía completa --- HMAC,
CMAC, CBC-MAC, las funciones hash aprobadas y las dos funciones de
derivación de SP~800-90A --- y da la fórmula de la entropía de una salida
acondicionada en función de la longitud de la entrada, la entropía de la
entrada y la anchura interna. La implementación de referencia de los
estimadores es pública~\cite{nist_ea_github}.

Las dos normas coinciden en lo que no pueden hacer las baterías
estadísticas más antiguas. SP~800-22~\cite{bassham_2010_sp80022}
comprueba la uniformidad y la independencia de un flujo de bits; un
generador determinista con entropía nula la supera por construcción, y
también una fuente física débil situada detrás de un postprocesado
criptográfico~\cite{saarinen_2021_entropy}. La decisión del NIST de 2022
de revisar la publicación lo dice expresamente: ``rechazando su uso para
evaluar generadores criptográficos de números aleatorios''~\cite{nist_2022_decision}.
Los datos empíricos lo confirman: un generador certificado de tarjeta
inteligente mostró sesgos persistentes aun superando la mayoría de los
tests~\cite{hurleysmith_2018_certifiably,
hurleysmith_2018_great}, y el generador del ESP32 supera esas mismas
baterías con la radio apagada, cuando la propia documentación de
Espressif dice que devuelve valores pseudoaleatorios~\cite{espressif_idf_random, espressif_trm_esp32,
blancoromero_2026_entropy}. La literatura académica sobre ese generador
solo aplica SP~800-22 o Dieharder~\cite{castillo_2026_automated,
hidayatulloh_2024_performance}; no se encontró ninguna caracterización
suya según SP~800-90B.

\section{Núcleos simétricos para hardware con recursos limitados}
\label{sec:soa_cores}

El motor necesita un cifrador de bloque y, si lo usa, una función hash.
Para AES-128 (diez rondas~\cite{nist_2023_fips197}), el espacio de diseño
va desde cauces segmentados totalmente desenrollados, orientados a la
tasa de bits, hasta rutas de datos de 8 bits, orientadas al área: Good y
Benaissa publican 124 slices y dos block RAM a 2,2~Mbit/s en una
Spartan-II en un extremo y 25~Gbit/s en el
otro~\cite{good_2005_fastest}; el diseño compacto de Chodowiec y Gaj usa
222 slices y tres block RAM~\cite{chodowiec_2003_compact}; y las
realizaciones ASIC más pequeñas llegan a 2,4~kGE con una ruta de datos de
8 bits a 226 ciclos por bloque~\cite{moradi2011pushing}, en la línea de
las S-box en cuerpo compuesto de Satoh~\cite{satoh_2001_compact},
Canright~\cite{canright_2005_compact} y
Boyar--Peralta~\cite{boyar_2010_minimization}. Las construcciones en
cuerpo compuesto optimizan el número de puertas en ASIC y pierden su
ventaja en una LUT de FPGA, donde una tabla directa es a la vez más
rápida y más sencilla. Existen diseños de referencia abiertos con cifras
publicadas para Artix-7: un cifrador de una ronda por ciclo con 1104~LUT
a 294~MHz~\cite{hadipour_aesvhdl}, y un cifrador/descifrador iterativo de
32 bits con 2298 slices a 97~MHz~\cite{strombergson_aes}. Para un motor
cuya salida sale por un enlace I2C, cualquiera de ellos es mucho más
rápido que el enlace, y la elección la dictan la verificabilidad y el
área, no la velocidad.

SHA-256~\cite{nist_2015_fips1804}, en su forma iterativa, cuesta unos 750
slices a 108~MHz en Artix-7~\cite{strombergson_sha256}, lo que es
comparable a un AES compacto; su papel en un generador sería el de
acondicionador aprobado de SP~800-90B. Los generadores deterministas de
SP~800-90A~\cite{barker_2015_sp80090ar1} --- Hash\_DRBG, HMAC\_DRBG y
CTR\_DRBG --- se diferencian sobre todo en el coste: CTR\_DRBG con
AES-128 mantiene un estado de 256~bits más un contador y produce
128~bits por cada cifrado de bloque, Hash\_DRBG necesita sumas modulares
de 440 bits, y HMAC\_DRBG, dos compresiones de hash por bloque de salida.
Sin función de derivación, CTR\_DRBG requiere una entrada con entropía
completa de exactamente la longitud de la semilla; con ella, que es la
configuración habitual en la práctica, acepta entradas arbitrarias. El
análisis de Woodage y Shumow advierte de que la flexibilidad de la norma
permite elecciones que empeoran la fuga del estado, y recomienda
actualizar tras cada operación de generación y resembrar muy por debajo
del límite de $2^{48}$~\cite{woodage_2019_analysis}; SP~800-90C,
definitiva desde 2025, establece el marco en el que se combinan una
fuente física y un DRBG~\cite{barker_2025_sp80090c}. Para la
autenticación, AES-CMAC (SP~800-38B~\cite{nist_2016_sp80038b},
RFC~4493~\cite{song_2006_rfc4493}) cuesta un cifrado de bloque por cada
bloque de mensaje y es en sí mismo un acondicionador aprobado. Los
vectores oficiales de todos ellos son
públicos~\cite{nist_cavp_blockciphers, nist_cavp_modes, nist_cavp_drbg,
nist_acvp_server}.

\section{Elementos seguros comerciales}
\label{sec:soa_se}

La categoría de producto a la que pertenece este motor existe. Un
elemento seguro es un pequeño circuito integrado que genera y almacena
claves y realiza operaciones simétricas y asimétricas para un
microcontrolador anfitrión, casi siempre por I2C. Cuatro componentes
dominan el segmento de bajo coste: el ATECC608B de Microchip, el SE050 de
NXP, el OPTIGA Trust~M de Infineon y el STSAFE-A110 de ST~\cite{microchip_atecc608b_ds,
nxp_se050_ds, infineon_optiga_trustm_ds, st_stsafe_a110_ds}. El
ATECC608B, a 0,90~\$ por unidad suelta~\cite{digikey_prices_2026}, ofrece
AES-128, SHA-256/HMAC, ECDSA/ECDH sobre P-256, dieciséis ranuras de clave
y una interfaz I2C, y su fuente de entropía cuenta con una validación del
NIST según SP~800-90B, obtenida precisamente con osciladores en
anillo~\cite{nist_esv_e46}. Dos de sus competidores no hacen ninguna
referencia a SP~800-90 en sus hojas de datos; el tercero afirma que su
generador está ``diseñado para cumplir'' la norma. La lección para este
trabajo es dónde puede competir un diseño académico y dónde no: no en
precio, donde un componente certificado cuesta menos de un euro, sino en
la transparencia de la entropía, que es justo lo que un componente de
caja negra no puede ofrecer y lo que las normas exigen ahora.

\section{De la FPGA al ASIC: qué cambia en la fuente de entropía}
\label{sec:soa_asic}

La lógica digital se traslada de una FPGA a celdas estándar con una
penalización conocida --- Kuon y Rose midieron 40$\times$ en área,
3,2$\times$ en retardo y 12$\times$ en potencia dinámica para diseños solo
de lógica en 90~nm, y 27--49$\times$ en área para AES en
concreto~\cite{kuon2006}; el oscilador en anillo no se traslada en
absoluto, porque en silicio se construye con celdas
inversoras cuyo ruido, variabilidad y acoplamiento difieren de los de una
cadena de LUT. Tres resultados publicados definen lo que debe añadir una
implementación en silicio. Benea \emph{et al.} muestran que el ruido de
parpadeo, cuyo peso crece en los nodos más recientes, cambia el carácter
del jitter y, contra lo esperado, que tratarlo como fuente legítima puede
cuadruplicar la tasa de bits a igual entropía~\cite{benea_2024_flicker}.
Yang, Blaauw y Sylvester muestran que la variación de proceso --- mayor en
40~nm que en 180~nm --- debe absorberse con un anillo configurable y un
lazo de ajuste para que el generador sobreviva entre $-40$ y
120~$^\circ$C y entre 0,6 y 0,9~V~\cite{yang_2016_jssc}. Y dos ataques
muestran por qué un dispositivo de seguridad no puede ignorar su entorno:
una frecuencia inyectada en la alimentación engancha los anillos y
eliminó el jitter de un microcontrolador seguro comercial, con lo que su
espacio de claves pasó de $2^{64}$ a 3300~\cite{markettos_2009_frequency},
y unos pocos cientos de microvatios de campo radiado sesgaron un
generador de cincuenta anillos hasta que falló todos los
tests~\cite{bayon_2012_contactless}; la defensa publicada combina el
filtrado de la alimentación con el mismo lazo de
calibración~\cite{yang_2016_jssc}, y el envejecimiento del dispositivo
por inestabilidad por polarización y temperatura y por inyección de
portadores calientes degrada con los años los generadores mal
temporizados, mientras deja intactos los bien temporizados.

El acceso al silicio es más abierto de lo que parece. La Universidad de
Sevilla y el IMSE-CNM son miembros de Europractice~\cite{europractice_memberlist},
cuya lista de precios académicos es pública: un bloque mini@sic cuesta
3\,430~\texteuro{} en UMC~180~nm y 3\,691~\texteuro{} en TSMC~65~nm, y el
SG13G2 de IHP se ofrece a 5\,110~\texteuro{}/mm$^2$, con una tirada
gratuita para diseños de código abierto de menos de 2~mm$^2$~\cite{europractice_prices2026, europractice_minisic,
ihp_openpdk}. El área de seguridad hardware del instituto ya ha llevado
primitivas de osciladores en anillo desde FPGA Xilinx de la serie 7,
evaluadas con SP~800-22 y SP~800-90B, hasta silicio de 65~nm
caracterizado~\cite{imse_area_hwsec,
martinez2022efficientropuf, ortega2024vlsi65nm, ortega2026ropuf65nm}, y
ha fabricado antes circuitos criptográficos~\cite{mora2020trivium}. Las
bibliotecas abiertas de celdas estándar que hacen posible un capítulo de
síntesis --- IHP SG13G2, SkyWater SKY130, GlobalFoundries GF180MCU y la
predictiva Nangate 45~nm --- vienen con vistas liberty, LEF y GDS, y la
literatura compara este tipo de diseños en puertas equivalentes
normalizadas a la NAND2 mínima de la biblioteca~\cite{ihp_openpdk_repo, skywater_pdk_cells, gf180mcu_pdk,
nangate45, kaps2019lwc}, con la advertencia de que el área de síntesis
lógica es una cota inferior que crece en un factor de 1,7 a 4,6 al pasar
al área de núcleo tras el rutado.

\section{Hueco que aborda este trabajo}
\label{sec:gap}

No se reivindica una novedad absoluta. El ERO y su modelo, la medida
embebida del jitter, el estudio comparativo de núcleos conformes a AIS,
el COSO calibrado en dispositivos de la serie 7 y una implementación
abierta sobre el mismo dado de la Basys~3 existen todos, y el propio
grupo del instituto ha pasado de las FPGA de la serie 7 al silicio de
65~nm con primitivas de osciladores en anillo. Lo que \emph{no se
encontró} en la literatura consultada, y lo que este trabajo se propone
aportar de forma modesta pero verificable, es lo siguiente.

\begin{enumerate}
\item Una cifra de jitter por periodo para la Artix-7 XC7A35T en
unidades absolutas, obtenida dentro del dispositivo con un método
diferencial, con las componentes térmica y de parpadeo separadas y con su
incertidumbre declarada, dado que los coeficientes publicados para
Artix-7 están en unidades del contador.
\item La cadena completa \emph{jitter medido $\to$ factor de calidad $\to$
cota de min-entropía $\to$ divisor $\to$ umbrales de los tests de salud},
aplicada a una plataforma y contrastada con los diez estimadores no
i.i.d.\ de SP~800-90B y con los requisitos PTG.2 de AIS~20/31
\emph{versión~3.0}, que la mayor parte de la literatura anterior no pudo
usar.
\item Una caracterización según SP~800-90B del generador del ESP32 en sus
estados de funcionamiento, y una comparación controlada, con la misma
metodología, entre una fuente de caja blanca en FPGA y una comercial de
caja negra.
\item El uso de un microcontrolador de uso común, con sus propios
aceleradores criptográficos, como oráculo independiente de un motor en
FPGA, sobre un bus real, con pruebas de respuesta conocida, de
interoperabilidad y de reto-respuesta.
\item Un motor de generación de claves de una sola primitiva --- AES para
el acondicionamiento, la generación y la autenticación --- escrito como
RTL portable, verificado contra vectores oficiales, sintetizado sobre una
biblioteca abierta de 130~nm con el resultado expresado como cota
inferior, y presupuestado para una lanzadera académica.
\end{enumerate}

Cada punto se formula como ``no encontrado'' y no como ``no existe''. Una
afirmación anterior de este tipo en las notas de trabajo --- que no se
había publicado ninguna medida del jitter en Artix-7 --- tuvo que
corregirse cuando se encontró el artículo de Benea \emph{et al.}, y la
corrección se mantiene visible allí.
